\section{Expected Results and Discussion Framework}

This section describes the expected outcomes of the study and the framework
for interpreting experimental findings.

The objective is not to assume that prompt engineering will completely replace
traditional algorithm training, but rather to investigate how AI-assisted
programming may change the skills required for software engineering.


\subsection{Expected Performance Trends}


The study expects that structured prompt strategies may improve LLM performance
compared with simple prompting.

Potential improvements may appear in:

\begin{itemize}

\item Solution correctness.

\item Algorithm selection.

\item Complexity optimization.

\item Explanation quality.

\end{itemize}


However, the improvement may vary depending on problem difficulty and
algorithm category.



\subsection{Possible Experimental Outcomes}


Several possible outcomes are considered.


\subsubsection{Outcome 1: Significant Improvement}


If optimized prompting produces substantially better results, this may suggest
that prompt engineering can effectively enhance AI-assisted algorithmic
problem solving.


This result would support the hypothesis that structured interaction with LLMs
is an important engineering skill.



\subsubsection{Outcome 2: Limited Improvement}


If optimized prompts only provide small improvements, this may indicate that
fundamental algorithm knowledge remains important even when using LLM tools.


This result would suggest that prompt engineering should be considered a
complementary skill rather than a replacement.



\subsubsection{Outcome 3: Performance Depends on Problem Type}


If improvements vary across categories, this may reveal that AI assistance is
more effective for certain algorithmic patterns.


For example:

\begin{itemize}

\item Simple data structure problems may receive limited benefits.

\item Complex reasoning problems may benefit more from structured prompts.

\end{itemize}



\subsection{Implications for Software Engineering Education}


The results may provide insights into how computer science education could
adapt to the LLM era.


Possible implications include:


\begin{itemize}

\item Teaching students how to collaborate with AI systems.

\item Combining algorithm fundamentals with prompt engineering skills.

\item Focusing less on memorization and more on verification and optimization.

\end{itemize}



\subsection{Limitations}


This research has several potential limitations:


\begin{itemize}

\item Results may depend on the selected LLM models.

\item Benchmark problems may not represent all software engineering tasks.

\item Prompt performance may change as models evolve.

\item Algorithm interviews evaluate human reasoning rather than only generated
solutions.

\end{itemize}



Future work may extend the benchmark by including more models, larger datasets,
and human-AI comparison studies.



\subsection{Discussion Direction}


The final discussion will focus on the following question:


\begin{quote}

Will future software engineers need fewer algorithm practice hours, or will
they need a new combination of algorithm knowledge and AI collaboration skills?

\end{quote}


The answer will be determined through empirical evaluation rather than
speculation.
